\section{Discussion}\label{sec:discussion}

Decomposing the error of DA schemes along the chain of Sect.~\ref{sec:chain}, for the mean
and for the reported uncertainty, changes how their performance can be read. With a perfect
model, the deficit of model-based DA with respect to the best schemes is mostly reducible:
on L96, the excess of the ensemble schemes over the best scheme splits into comparable shares
of restriction, due to filtering, approximation, due to the ensemble size, and an unattributed
term (Sect.~\ref{sec:q1}). The first two shrink with smoothing and with more members. The third
includes a structural misspecification of the posterior family --- Gaussian updates, the mode
instead of the mean, a hard constraint --- which dominates the budget of the 4D-Var schemes, so
that misspecification is not zero even when the dynamical model is exact. This deficit is not
calibrated either: the DA ensembles are over-confident, and the more so as their mean improves
(Sect.~\ref{sec:q1-uncertainty}). Smoothing and larger ensembles thus reduce the error of the
mean, not the miscalibration of the reported uncertainty.

Under model error, misspecification dominates the budget of every model-based scheme. On L96, it
multiplies their normalised MSE by two to four, and the contrasts that reduced the perfect-model
excess no longer help (Sect.~\ref{sec:q2}). Representing the uncertainty of the model recovers
part of the loss --- a weak constraint reduces the error of the hard one by a third to two
fifths --- without closing the gap. The flow-matching score shows that misspecification
degrades the reported distribution even more than the mean: the DA ensembles become more
over-confident and short-tailed, and their score at large $\tau$ falls below that of
climatology, whereas the learned schemes report a distribution whose skill and calibration
are unchanged. The advantage of the learned schemes is thus larger for the distribution than
for the mean, and the large-$\tau$ score reveals misspecification where the mean-based scores
only suggest it.

The comparison of priors at matched likelihood (Sect.~\ref{sec:q3}) qualifies the common
reading that learned schemes win because they learn a better prior. A learned prior
outperforms the exact equations used by a 30-member smoother or by a hard constraint, but not
by a 100-member smoother: the ordering is set by the approximation budget of the scheme that
uses the prior, not by the fidelity of the prior. Conditioning a learned prior on the model
parameters brings little with a perfect model, and becomes a misspecification channel under
model error unless the prior is trained with noisy parameters. Composing an amortised mean
with a guided sampler gives the best mean in both case studies, but an over-confident
distribution, and the composition is what makes the learned prior usable on QG.

Model error also changes what the observations bring (Sect.~\ref{sec:q1-density}). Under
model error, additional or more accurate observations bring little to the DA schemes, whose
error is set by their model, whereas the learned schemes keep their perfect-model response to
the observing system. DA-based assessments of an observing system, such as observing system
simulation experiments, may therefore underestimate its value, and the more so as the model is
imperfect. Classical DA nevertheless keeps clear strengths. It is better when observations are
very sparse, and, together with SDA, it generalises beyond the observing systems seen in
training, because it keeps the observation operator and the error covariance explicit: the
learned schemes degrade when a window holds fewer or many more observations than in training,
and transfer less well between observing systems. On QG, with a realistic observing system,
strong-constraint 4D-Var gives the best mean with a perfect model. DA also needs no training
set, whereas the learned schemes we compare are trained on simulations of the true system.
\todo{Back to the puzzle: what the budgets suggest for reanalyses versus observation-only
products, stated as hypotheses for real systems.}

These conclusions come with limitations. The irreducible term $\mathrm{mmse}_p(\tau)$ is
unknown, so that every component is a share of an excess over the best scheme, and the cross
term between misspecification and approximation is bounded rather than identified. The L96
imperfect-model experiment combines a parametric error with a structural one, as the DA model
is truncated to the observed variables \needsrun{\So-full rows}, and the learned rows of this
experiment are oracle bounds, since the learned schemes are trained on the true system
\needsrun{biased-model training arms}. The flow-matching score is evaluated per element, with
sliced evaluations as a check of the dependence structure, and its shape diagnostic has a small
negative floor for finite and over-confident ensembles. Finally, the cross-class conclusions
rest mostly on L96: on QG, only D-Unet and D-Unet+SDA are available among the learned schemes,
and the L63 case study is to be completed.

This study relates to the literature on model error in DA --- weak-constraint
4D-Var~\cite{tremolet_accounting_2006}, inflation and model-error
estimation~\cite{carrassi_data_2018,farchi_using_2021}, assimilation in the unstable
subspace~\cite{trevisan_assimilation_2004,bocquet_four-dimensional_2017} --- and on learned and
hybrid DA~\cite{brajard_combining_2021,bocquet_bayesian_2020,fablet_learning_2021}, including
score-based and flow-based DA~\cite{rozet_score-based_2023}. \todo{Proper scores and their
decompositions; mismatched estimation.} To our knowledge, no controlled study has measured
\emph{where} the error of a DA scheme comes from --- restriction, misspecification or
approximation --- for the mean and for the reported uncertainty, at matched information across
classical and learned schemes.

Several questions remain open. The first concerns the nature of model error: separating
parametric from structural error, as the full-model variant of L96 begins to do, and varying the
bias with respect to the stated uncertainty, which is tractable on L63, would show when
representing model uncertainty converts misspecification back into restriction. The second
concerns the training data: training the learned schemes on the biased model, or on reanalyses,
would price the value of simulations of the true system, which our learned schemes use and
classical DA does not, and would tell whether their robustness to model error survives without
it. The third concerns the scores: joint or sliced versions of the flow-matching score would
assess the dependence structure of the reported uncertainty more completely, and its
calibration and shape diagnostics could guide the post-processing of ensembles. More broadly,
the operator~\eqref{eq:psi} suggests writing every DA scheme as a flow and composing the mean of
one with the anomaly of another, as D-Unet+SDA does, with a calibration step to correct the
over-confidence of the composition. Finally, the framework extends to the joint estimation of
the state and of the model parameters, by including $\thv$ in the target of the flow, which
would bring the misspecification of the parameters into the same budget.
